\mxnewpage\section*{EV3-T1 -- Blanc 2}\label{EV3-T1}
\addcontentsline{toc}{section}{EV3-T1 -- Blanc 2}
\textbf{Sujet d’entraînement}. Programme étudié des chapitres 1 à 10.\par Durée : 2 h. Sans calculatrice. Note sur 20. 1–3 avril 2027.
\subsection*{Partie 1 -- Automatismes (6 points)}
Douze questions, \num{0,5} point chacune. Une seule réponse correcte ; aucune justification demandée. Une réponse incorrecte ou absente rapporte 0, sans retrait de points.
\begin{enumerate}
\item Calculer $6+2\times3$.\par A. $12$\quad B. $24$\quad C. $11$\quad D. $15$
\item Calculer $\frac12+\frac13$.\par A. $\frac{2}{5}$\quad B. $\frac{5}{3}$\quad C. $\frac{5}{6}$\quad D. $\frac{1}{6}$
\item Calculer $\dfrac{2^6}{2^{4}}$.\par A. $1$\quad B. $11$\quad C. $2$\quad D. $4$
\item Une valeur de 50 augmente de \pc{20}. Valeur finale ?\par A. $10$\quad B. $40$\quad C. $70$\quad D. $60$
\item Une valeur de 80 est obtenue après une baisse de \pc{20}. Valeur initiale ?\par A. $64$\quad B. $96$\quad C. $100$\quad D. $60$
\item Hausse de \pc{25} puis baisse de \pc{20} : coefficient global ?\par A. $\frac{21}{20}$\quad B. $\frac{9}{10}$\quad C. $1$\quad D. $\frac{4}{5}$
\item La solution de $2x+6=0$ vaut\par A. $12$\quad B. $-12$\quad C. $-3$\quad D. $3$
\item La pente entre $\pt{0;6}$ et $\pt{2;12}$ vaut\par A. $-3$\quad B. $3$\quad C. $2$\quad D. $6$
\item La moyenne de $6, 8, 10$ vaut\par A. $9$\quad B. $8$\quad C. $24$\quad D. $6$
\item Un secteur de $90^\circ$ représente quelle fraction du total ?\par A. $\frac{1}{2}$\quad B. $\frac{1}{3}$\quad C. $\frac{1}{4}$\quad D. $\frac{3}{4}$
\item Une probabilité vaut $\dfrac{2}{5}$. Son complément vaut\par A. $\frac{5}{2}$\quad B. $\frac{2}{5}$\quad C. $\frac{3}{5}$\quad D. $\frac{1}{5}$
\item Pour $f(x)=x^2+6$, $f(-2)$ vaut\par A. $10$\quad B. $8$\quad C. $24$\quad D. $2$
\end{enumerate}
\mxnewpage\subsection*{Partie 2 -- Exercices indépendants (14 points)}
Toute démarche doit être justifiée. Les résultats exacts sont attendus sauf indication contraire.
\mxneed{12}
\subsection*{Exercice 1 -- Un gain composé (5 points)}
Un jeu consiste à effectuer deux essais indépendants, avec une probabilité de succès de $\dfrac12$ à chaque essai. Chaque succès rapporte 3 unités monétaires ; la participation coûte 2 unités. On note $X$ le gain net, c’est-à-dire le total reçu moins la participation payée.
\begin{enumerate}
\item Établir la loi de probabilité du gain net $X$. \hfill(1 pt)
\item Calculer son espérance et sa variance. \hfill(2 pts)
\item Déterminer la participation qui rendrait le jeu équitable. Expliquer ce que deviendrait l’écart type du gain net. \hfill(2 pts)
\end{enumerate}
\mxneed{12}
\subsection*{Exercice 2 -- Pic d’un modèle (5 points)}
Une grandeur est modélisée pour $t\geq0$ par la fonction $f(t)=(t+2)e^{-t}$. Le temps $t$ est exprimé dans une unité fixée. On cherche sa valeur maximale, puis une approximation au voisinage de l’instant initial.
\begin{enumerate}
\item Dériver f. \hfill(1 pt)
\item Étudier les variations et donner le maximum. \hfill(2 pts)
\item Écrire la tangente en 0 et approcher $f(\num{0,1})$. \hfill(2 pts)
\end{enumerate}
\mxneed{12}
\subsection*{Exercice 3 -- Cercle trigonométrique (4 points)}
Le plan est muni d’un repère orthonormé. Sur le cercle trigonométrique de centre $O$, les points $A$ et $B$ sont les images respectives des réels $\dfrac\pi6$ et $\dfrac{5\pi}6$.
\begin{enumerate}
\item Donner leurs coordonnées. \hfill(1 pt)
\item Calculer AB. \hfill(1 pt)
\item Donner l’équation de leur médiatrice et justifier géométriquement. \hfill(2 pts)
\end{enumerate}
